\begin{center}
\textbf{\large{Abstract}	}
\end{center}

\addcontentsline{toc}{chapter}{Abstract}
\begin{small}

\textbf{\textit{Abstract: }}Bacteriophages are viruses that infect bacteria and are regarded as a potential alternative to antibiotics through phage therapy. In therapeutic applications, accurately identifying the lifestyle of a phage, virulent or temperate, is an important requirement, since most therapeutic applications prioritize virulent phages; temperate phages may integrate their genetic material into the host genome through the lysogenic cycle and contribute to the spread of virulence genes or antibiotic resistance genes via horizontal gene transfer, thereby increasing therapeutic risk. Traditional culture-based laboratory methods are typically time- and resource-intensive and are difficult to apply to the majority of phages that remain uncultured or unlabeled, motivating the need for computational classification methods. Existing tools such as DeePhage, PhaTYP, and ProkBERT achieve promising performance but still have their own limitations.

This thesis proposes PhaBERT, a deep learning architecture that combines the DNABERT-2 genome foundation model with a multi-kernel convolution (MKC) module. Through this design, PhaBERT leverages pre-trained knowledge from multi-species genomic data while operating directly on raw DNA sequences through BPE tokenization without information leakage.

Experiments on four contig length groups (100-400 bp, 400-800 bp, 800-1,200 bp, 1,200-1,800 bp) with 5-fold stratified cross-validation show that PhaBERT achieves accuracies of 82.01\%, 87.33\%, 89.90\%, and 91.38\% respectively, improving by 0.36--6.65\% over state-of-the-art methods including PhaTYP, DeePhage, and ProkBERT. An ablation study indicates that the MKC module contributes accuracy improvements ranging from 0.49\% to 3.02\% over the DNABERT-2 baseline using only a linear classification head, with the improvement increasing as contig length grows. The model maintains a balance between sensitivity (82.07\%--92.38\%) and specificity (81.86\%--87.63\%) across the length ranges examined. These results suggest the effectiveness of combining a genome foundation model with a task-specific architecture. This is a promising approach for biological sequence analysis tasks, especially metagenomic analysis involving incomplete or fragmented sequences.

\vspace*{1cm}
\textbf{\textit{Keywords}}: \textit{Bacteriophage}, \textit{Lifestyle classification}, \textit{Deep learning}, \textit{BERT}, \textit{Convolutional neural networks}, \textit{Bioinformatics}, \textit{Metagenomics}
\end{small}
